% !TEX root = ../../mixed-sklyanin-cubics.tex
\section{非 PI 変形の相対対角と一般点核の整延長}
\label{mix:sec:kernel-specialization}

定理\ref{mix:thm:h-exact-candidate}の一般点核を $h=1$ で調べるためには、
形式的な係数の代入ではなく、同じ代数族の上に proper な圏と perfect な核を置く必要がある。
本節では、次数付き ambient の有限 Koszul 解消と、射影積上の通常の coherent 対角を使う。
族全体が有限中心を持つことは仮定しない。
得られるのは、$h=1$ を含む実際の相対 smooth proper 模型と、
指定された一般点核を延長する perfect 整核である。
$h=1$ の通常の Ku 同値の存在は、
定理\ref{cmo:thm:fixed-zero-commutative-kernel}の直接構成で得られる。
本節の整核について残るのは、選択済みの一般点核との整合性を保ったまま、
特殊ファイバーでも unit/counit が同型となる延長を選べるかという問題である。

\subsection{局所および半局所の七線束模型}

\[
 R_1=\C[h]_{(h-1)},\qquad
 R_{01}=S^{-1}\C[h],\qquad
 S=\{f\in\C[h]\mid f(0)f(1)\ne0\}
\]
と置く。以下で $R$ はこのいずれかとする。
その Jacobson radical の生成元を
\[
 \varpi=h-1\quad(R=R_1),\qquad
 \varpi=h(h-1)\quad(R=R_{01})
\]
と書く。$R_{01}/(\varpi)=\C\times\C$ の二因子は $h=0,1$ に対応する。
$R$ は正則な半局所 PID で、$h^6+27$ は単元である。
式\eqref{mix:eq:h-mixed-family}の base change を
\[
 B=\mathscr B\otimes_{\C[h]}R,\qquad A=B/(w_h)
\]
と書く。

\begin{proposition}[非 PI 族の proper な七線束模型]
\label{mix:prop:h-semilocal-model}
$B,A$ および $B^{\mathrm{op}}\otimes_RB$、
$A^{\mathrm{op}}\otimes_RA$ は左右 noetherian である。
$B$ は $R$ に相対的な長さ七の augmentation Koszul 解消と
長さ七の双加群解消を持ち、その各項の階数は $\binom7j$ である。
これらは次数別に $R$-分裂し、通常の底変更を保つ。
相対 Gorenstein 型は
\[
 R\underline{\operatorname{Hom}}_B(R,B)\simeq R(7)[-7],
 \qquad
 R\underline{\operatorname{Hom}}_A(R,A)\simeq R(4)[-6]
\]
である。

$\mathcal O(i)=\pi A(i)$ とし、
\[
 E=\bigoplus_{i=0}^{6}\mathcal O(i),\qquad
 \Sigma_R=R\operatorname{End}(E)
\]
と置く。$\Sigma_R$ は $R$-proper であり、
\begin{equation}\label{mix:eq:h-local-line-hom}
 \operatorname{Ext}^p(\mathcal O(i),\mathcal O(j))=
 \begin{cases}
 A_{j-i},&p=0,\\
 (A_{i-j-4})^\vee,&p=5,\\
 0,&\text{その他}
 \end{cases}
\end{equation}
となる。各群は有限自由 $R$-加群である。
この dg 模型と例外列
$\langle\mathcal O,\mathcal O(1),\mathcal O(2),\mathcal O(3)\rangle$
による射影 $\operatorname{pr}$ に対し、
\[
 G_R=\operatorname{pr}\mathcal O(4)\oplus
     \operatorname{pr}\mathcal O(5)\oplus
     \operatorname{pr}\mathcal O(6),\qquad
 \Gamma_R=R\operatorname{End}(G_R)
\]
を残余模型とする。
両模型は剰余体および一般点への導来底変更と両立する。
両模型には $u\mapsto-u$ の相対作用がある。
\end{proposition}

\begin{proof}
命題\ref{mix:prop:h-flat}により、全代数は次数別有限自由である。
$h=0$ の ambient は七変数多項式環であり、
$h=1$ の ambient の Koszul 性と自己双対性は
命題\ref{cmo:prop:proper-contraction}の証明を $t=1$ で特殊化して得られる。
後者とその商は、系\ref{cmo:cor:relative-finite-center}によって有限中心上有限である。
従って $B/(\varpi)$ とその包絡環、および対応する $A$ の商は左右 noetherian である。

補題\paperref{main}{lp:lem:noetherian}の証明を $(\varpi)$ に適用する。
同補題で使う条件は、係数環の局所性そのものではなく、
$(\varpi)\subset\operatorname{Jac}(R)$ と、各内部次数での Krull 交叉である。
実際、
\[
 (B/(\varpi))[T]\twoheadrightarrow\operatorname{gr}_{(\varpi)}B
\]
があり、斉次イデアルの初項の有限生成元を持ち上げる。
固定内部次数で商は有限 $R$-加群なので、
$\bigcap_q\varpi^qM=0$ と Nakayama による同じ証明が働く。
包絡環には二つの内部次数の和を使う。
この方法は、一般ファイバーの PI 性を必要としない。

補題\paperref{main}{lp:lem:resolution}も同様に半局所環上で使える。
二つの剰余ファイバーの Koszul 空間は同じ階数と shift を持ち、
$R$ 上の有限射影加群は自由である。
augmentation と対角の有限自由解消を次数別に持ち上げると、
各段階の核は次数別有限自由で、底変更後の核も正しく計算される。
最後の核は剰余が零なので零である。
双対解消の単元成分を消去すれば上の ambient Gorenstein 式を得る。
正則中心三次 $w_h$ の二項解消から商の式が従う。

以上は命題\ref{cmo:prop:proper-contraction}の線束 Hom の計算と同じ入力である。
特に、補題\paperref{main}{lp:lem:local-cohomology}の次数 filtration の証明と
命題\paperref{main}{lp:prop:basechange}の底変更の証明が適用でき、
式\eqref{mix:eq:h-local-line-hom}を得る。
例外射影は有限個の evaluation cone であり、その Hom は $R$-perfect なので、
残余模型も proper で底変更を保つ。
対合は元の代数、線束、evaluation を保つ。
\end{proof}

\subsection{射影積における有限性}

通常の coherent 対角を使うため、射影積の模型を明記する。
$A$ と、同じ条件を満たすもう一つの中心三次商 $A'$ に対し、
\[
 C=A^{\mathrm{op}}\otimes_RA'
\]
を二重次数付き代数とする。
$\operatorname{QGr}^{(2)}C$ は、
各元が $C_{\ge n,\ge n}$ で消える加群を除いた商であり、
$\operatorname{qgr}^{(2)}C$ はその noetherian 対象の部分圏とする。
この積と二重次数対角は
Ballard--Farman \cite[Definitions~3.16, 3.18, Section~3.4]{CMOBallardFarman}
のものを使う。
本節では体上の結果をそのまま係数環 $R$ に適用せず、
以下の相対有限性を確認する。
必要な $A'$ は $A$ 自身、または
$R[X_1,\ldots,X_7]/(W_h^{\mathrm{cand}})$ である。
後者の $h=0,1$ のファイバーは滑らかである。

\begin{lemma}[ambient 解消による射影積の制御]
\label{mix:lem:h-projective-product}
上の二重次数射影圏では、次が成立する。
\begin{enumerate}
 \item 全ての線束
 $\mathcal O(i,j)=\pi(A^{\mathrm{op}}(i)\otimes_RA'(j))$
 は compact であり、それらは非有界導来圏を生成する。
 任意の abelian 圏の対象 $N\in\operatorname{QGr}^{(2)}C$ に対して
 $\operatorname{Ext}^q(\mathcal O(i,j),N)=0$ は $q>12$ で成立する。
 \item coherent な $N$ に対し
 $R\operatorname{Hom}(\mathcal O(i,j),N)$ は $R$-perfect である。
 $N$ が次数別 $R$-flat な加群で表示される場合、この計算は剰余体と一般点への底変更を保つ。
 \item coherent な $M,N$ に対し、各
 $\operatorname{Ext}^q(M,N)$ は有限 $R$-加群である。
\end{enumerate}
一変数の射影圏にも同じ結論が成立し、第一項の上界は六でよい。
\end{lemma}

\begin{proof}
まず中心三次商 $A=B/(w)$ の torsion を ambient $B$ で計算できることを確認する。
この一段の Tor--Hom 計算は左加群で記し、右加群には反対環で同じ計算を適用する。
$B_{<n}=B/B_{\ge n}$ とすると、$A$ の二項 $B$-自由解消より
\[
 \operatorname{Tor}_i^B(A,B_{<n})=0\quad(i>1).
\]
$\operatorname{Tor}_1$ は
$w:B_{<n}(-3)\to B_{<n}$ の核である。
$w$ の正則性により、この核の内部次数は $n,n+1,n+2$ に限られる。
従って $n'\ge n+3$ の遷移
\[
 \operatorname{Tor}_1^B(A,B_{<n'})
 \longrightarrow\operatorname{Tor}_1^B(A,B_{<n})
\]
は零である。
$I$ を graded $A$-injective とし、導来随伴を用いれば
\[
 \operatorname{Ext}_B^q(B_{<n},I)
 =\operatorname{Hom}_A
       \bigl(\operatorname{Tor}_q^B(A,B_{<n}),I\bigr).
\]
右辺の正次数部分は $n$ に関する直極限で零になる。
従って $I$ は $B$-torsion functor に関して acyclic であり、
$A$-加群上で
\begin{equation}\label{mix:eq:h-torsion-independence}
 R\tau_A=R\tau_B,\qquad RQ_A=RQ_B
\end{equation}
が成立する。ここで $RQ$ は導来 saturation functor である。

$B$ の相対対角解消は長さ七であり、$B_{<n}$ は $R$-free なので
$\operatorname{pd}_B B_{<n}\le7$ である。
noetherian 性から、この加群には有限階数の項を持つ有限自由解消を選べる。
従って $R\tau_B$ の cohomological dimension は高々七で、
coproduct と可換である。$RQ_B$ の正次数は六以下になる。
式\eqref{mix:eq:h-torsion-independence}により同じことが $A$ に成立する。

二重次数の場合、二つの saturation の合成が積の saturation を計算する。
これは \cite[Proposition~3.41]{CMOBallardFarman} の truncation による計算である。
相対の場合も、各因子が $R$-flat なので加群の restriction は injective を保ち、
必要な truncation の有限自由解消と coproduct の可換性は直前に確認した。
tensor を導来 tensor として同じ計算を行うと
\[
 RQ_C\simeq RQ_A\,RQ_{A'}
           \simeq RQ_B\,RQ_{B'}
\]
を得る。上界は $6+6=12$ で、coproduct とも可換である。
内部次数 $(i,j)$ の評価は線束からの導来 Hom だから、これらの線束は compact である。
全線束からの Hom が零なら導来 saturation の全次数が零となるので、生成性も従う。

次に $D=B^{\mathrm{op}}\otimes_RB'$ と置く。
その相対対角解消は長さ十四であり、
$\operatorname{gldim}R=1$ から $\operatorname{gldim}D\le15$ である。
$D$ は noetherian で、coherent な二重次数 $C$-加群は有限生成 $D$-加群でもある。
従って有限階数の二重次数自由加群による有限解消を持つ。
最後の有限射影加群が自由であることは、
正次数イデアルで割った有限射影 $R$-加群の基底を持ち上げれば分かる。
この解消を tails へ移す。
ambient 線束の cohomology は相対 Gorenstein 計算により各因子で次数零と六にのみあり、
各群は有限自由 $R$-加群である。
二つの有限解消を tensor した K\"unneth 計算も同様である。
従って $N$ の線束からの Hom は有界で有限 $R$-cohomology を持つ。
$R$ が正則なのでこれは $R$-perfect である。
$N$ が次数別 $R$-flat なら、解消の各 syzygy も次数別 $R$-flat であり、
各内部次数で分裂した解消を選べる。
ambient 線束の底変更と合わせると第二項の底変更も従う。

最後に、$M$ を有限線束和による上に有界な解消で計算する。
$C$ の noetherian 性により各項を有限和にできる。
その hyper-Ext spectral sequence では、
固定した全次数に寄与する項は有限個であり、
第二項により各項が有限 $R$-加群である。
これで第三項を得る。
\end{proof}

\subsection{coherent な平坦対象の compact 性の持上げ}

次の補題は、射影積上の通常の加群を使う箇所である。
proper な dg 代数のファイバーが smooth であるという条件だけからの推論ではない。

\begin{lemma}[有限 Ext と普通の平坦対象による持上げ]
\label{mix:lem:coherent-compact-lift}
補題\ref{mix:lem:h-projective-product}の一変数または二変数射影圏を $\mathcal C_R$ とする。
coherent な $M\in\mathcal C_R$ に対して
$\varpi:M\to M$ が単射であり、
各閉ファイバー $M_s$ が $D(\mathcal C_s)$ で compact であるとする。
このとき $M$ は $D(\mathcal C_R)$ で compact である。
\end{lemma}

\begin{proof}
閉ファイバーの線束は compact generator の集合なので、
$M_s$ は有限個の線束の有限な cone、shift、直和、直和因子から得られる。
線束の cohomological dimension の一様な上界を使うと、
ある $d$ があり、全ての coherent な $N_s$ について
\begin{equation}\label{mix:eq:h-special-ext-bound}
 \operatorname{Ext}^q_{\mathcal C_s}(M_s,N_s)=0
 \qquad(q>d)
\end{equation}
となる。閉点は一つまたは二つしかないので、$d$ はその最大値に取れる。

まず $\varpi:N\to N$ も単射とする。
特殊ファイバーの包含を $i$ と書く。
$M$ の $R$-flatness による導来随伴から
\[
 R\operatorname{Hom}_{\mathcal C_R}(M,i_*N_0)
 \simeq R\operatorname{Hom}_{\mathcal C_{R/(\varpi)}}(M_0,N_0)
\]
を得る。ここでは $R_{01}/(\varpi)=\C\times\C$ の場合は二つの Hom の直積である。
短完全列
$0\to N\xrightarrow{\varpi}N\to i_*N_0\to0$
と式\eqref{mix:eq:h-special-ext-bound}より、$q>d$ で
\[
 \varpi:\operatorname{Ext}^q(M,N)\longrightarrow
                   \operatorname{Ext}^q(M,N)
\]
は全射となる。
補題\ref{mix:lem:h-projective-product}の有限性と
$(\varpi)\subset\operatorname{Jac}(R)$ による Nakayama から、
この Ext 群は零である。

一般の coherent な $N$ に対しては、有限線束和からの全射
$P\twoheadrightarrow N$ を取り、核を $K$ とする。
$K\subset P$ なので $\varpi$ は $K$ 上単射である。
直前の消滅を $P,K$ に適用し、長完全列を使えば、
同じ $q>d$ で $\operatorname{Ext}^q(M,N)=0$ を得る。
従ってこの議論は、$N[\varpi]\ne0$ の対象も含む。

$M$ に有限線束和から順に全射を取り、$n>d$ に対して
\[
 0\longrightarrow K_n\longrightarrow P_{n-1}\longrightarrow\cdots
 \longrightarrow P_0\longrightarrow M\longrightarrow0
\]
を作る。各 syzygy は coherent である。
この列の境界射 $M\to K_n[n]$ は
$\operatorname{Ext}^n(M,K_n)=0$ により零である。
従って $M$ は有限複体
$[P_{n-1}\to\cdots\to P_0]$ の直和因子である。
各 $P_i$ は compact なので、$M$ も compact である。
この証明は coherent 圏に projective object があることを仮定しない。
\end{proof}

\begin{theorem}[零点と一点を含む実際の相対 perfect 対角]
\label{mix:thm:h-semilocal-smooth}
命題\ref{mix:prop:h-semilocal-model}の $\Sigma_R,\Gamma_R$ は相対 smooth proper である。
また
\[
 D^b(\operatorname{qgr}A)=\operatorname{Perf}(\Sigma_R)
\]
であり、全 coherent Kuznetsov 成分は $\operatorname{Perf}(\Gamma_R)$ である。
これらの等式は dg Morita 同値を意味する。
特に $R=R_{01}$ は $h=0,1$ を同じ相対 smooth proper 模型に含め、
その模型は対合を保つ。
\end{theorem}

\begin{proof}
ambient augmentation 解消を $A$ へ tensor した複体の cohomology は
$R,R(-3)$ だけであり、どちらも torsion である。
従って tails でこの長さ七の複体は完全となる。
両端の線束を順に解けば、七つの連続した線束が全 twist を厚生成する。
補題\ref{mix:lem:h-projective-product}の一変数版により、
$E$ は $D(\operatorname{QGr}A)$ の compact generator である。

二重次数の対角を
\[
 (\Delta_A)_{i,j}=A_{i+j}\qquad(i,j\in\mathbb Z)
\]
とする。射影圏では $i,j\ge0$ の部分 $\Delta_A^+$ で表示できる。
商 $\Delta_A/\Delta_A^+$ の各元は十分大きい二重次数で消えるからである。
$A$ は次数一で生成されるので、$\Delta_A^+$ は $(0,0)$ の元 $1$ で生成される。
従って
\[
 \mathcal D_A=\pi\Delta_A^+
\]
は普通の coherent 対象であり、次数別に $R$-free である。
この具体的な対象の compact 性を証明する。

相対 Gorenstein 双対は線束 $\mathcal O(i)$ を左線束 $\mathcal O(-i)$ に移す。
torsion 対象の双対も torsion なので、この双対は線束の厚生成圏へ降りる。
$E^\vee=\bigoplus_{i=0}^{6}\mathcal O_{A^{\mathrm{op}}}(-i)$、
$F=E^\vee\boxtimes_RE$ と置くと、$F$ は積の compact generator である。
線束の K\"unneth と対角の評価により
\begin{equation}\label{mix:eq:h-projective-diagonal}
 \begin{split}
 R\operatorname{End}(F)&\simeq
       \Sigma_R^{\mathrm{op}}\otimes_R^{\mathbf L}\Sigma_R,\\
 R\operatorname{Hom}(F,\mathcal D_A)&\simeq\Sigma_R
 \end{split}
\end{equation}
を得る。第二式は左右作用を持つ式である。

この相対計算に使う条件を補足する。
graded 加群上の tensor--Hom evaluation を二つの saturation に通す計算は、
\cite[Propositions~3.42, 4.14, Theorem~4.15, Remark~4.17]{CMOBallardFarman}
の対角と外積の計算である。
ここでは、左右の相対 Gorenstein 型から得られる $\chi$、
式\eqref{mix:eq:h-torsion-independence}による有限 cohomological dimension、
次数別自由性を使って同じ導来 tensor 計算を行う。
各 truncation には有限階数自由解消を使い、基底環上の tensor は導来 tensor とする。
従って体上の tensor の完全性を暗黙に使う必要はない。
また、得られる evaluation と K\"unneth の写像は、
補題\ref{mix:lem:h-projective-product}によって両辺 $R$-perfect で底変更を保つ。
閉ファイバーでは上記文献の恒等式であり、
その cone の剰余が零であることと半局所 Nakayama からも
式\eqref{mix:eq:h-projective-diagonal}を確認できる。

$h=0$ の全 coherent 圏は滑らかな可換三次五様体の圏である。
$h=1$ の全 coherent 圏の smooth proper 性は
定理\ref{cmo:thm:exact-sixty-four}と $P_{64}(0)\ne0$ による。
従って式\eqref{mix:eq:h-projective-diagonal}の閉ファイバーで、
$\mathcal D_{A_s}$ は compact である。
補題\ref{mix:lem:coherent-compact-lift}を普通の平坦対象 $\mathcal D_A$ に適用すると、
$\mathcal D_A$ は compact となる。
compact generator による dg Morita 理論
\cite[Theorem~3.8(c)]{NFMKellerDg}と式\eqref{mix:eq:h-projective-diagonal}から
\[
 \Sigma_R\in\operatorname{Perf}
       (\Sigma_R^{\mathrm{op}}\otimes_R^{\mathbf L}\Sigma_R)
\]
を得る。有限例外列の admissible な直交成分にも smooth proper 性が移るので、
$\Gamma_R$ についても同じ結論となる。

全 coherent 圏との一致も確認する。
coherent な平坦対象 $M$ の閉ファイバーは、既知の smooth proper な全 coherent 圏の対象なので compact である。
従って同じ持上げ補題により $M$ は compact である。
一般の coherent 対象は有限線束和 $P$ の商として表示し、
平坦な核 $K\subset P$ を使えば compact となる。
逆に、有限線束複体とその直和因子は有界 coherent cohomology を持つ。
これで全 coherent 圏と compact 圏の一致を得る。
全構成の上で $u\mapsto-u$ は定義され、対角の同変構造も保たれる。
\end{proof}

同じ証明は、$R_1$ の有限体拡大における半局所整閉包 $\widetilde R$ と、
その各閉点での DVR $R'$ にも成立する。
$\widetilde R$ は有限平坦な半局所 Dedekind 環であり、
その Jacobson radical も単項で、剰余環は有限個の $\C$ の直積である。
各次数の自由性、特殊ファイバーの Koszul 型、
特殊ファイバーでの smooth proper 性が変わらないので、
上の $(\varpi)$ による議論をそのまま適用できる。
また $A$ と可換候補の混合射影積にも適用できる。
その閉ファイバーの積 dg 代数は smooth proper であり、
任意の coherent 対象は、線束からの Hom が proper であることから compact となる。
補題\ref{mix:lem:coherent-compact-lift}により、
相対混合積でも全ての coherent 対象が compact である。

\subsection{指定された一般点核を延長する}

\begin{theorem}[指定候補への perfect 整核の存在]
\label{mix:thm:h-integral-kernel}
$K=\C(h)$ とし、定理\ref{mix:thm:h-exact-candidate}で得た有限拡大 $K'/K$ と
指定候補への一般点比較核 $P_{K'}$ を固定する。
$R_1$ の $K'$ における整閉包を $\widetilde R$ とする。
これは $R_1$ 上有限平坦な半局所 Dedekind 環である。
その各閉点での局所化を $R'$ とする。
非可換側の残余模型を $\Gamma_{\widetilde R}$、
可換候補の残余模型を $\Lambda_{\widetilde R}$ とすると、
\[
 P_{\widetilde R}\in\operatorname{Perf}
  (\Gamma_{\widetilde R}^{\mathrm{op}}
       \otimes_{\widetilde R}^{\mathbf L}\Lambda_{\widetilde R}),
 \qquad P_{\widetilde R}\otimes_{\widetilde R}K'\simeq P_{K'}
\]
を満たす核が存在する。特に各 $R'$ 上で
\begin{equation}\label{mix:eq:h-integral-kernel}
 P_{R'}\in\operatorname{Perf}
       (\Gamma_{R'}^{\mathrm{op}}\otimes_{R'}^{\mathbf L}\Lambda_{R'}),
 \qquad
 P_{R'}\otimes_{R'}K'\simeq P_{K'}
\end{equation}
を満たす核が存在する。
従って両側に随伴を持つ実際の関手
\[
 \Phi_{P_{R'}}:\operatorname{Perf}(\Gamma_{R'})
       \longrightarrow\operatorname{Perf}(\Lambda_{R'})
\]
を構成でき、その一般ファイバーは指定された Ku 同値である。
この主張は $\Phi_{P_{R'}}$ の特殊ファイバーが同値であることを含まない。
\end{theorem}

\begin{proof}
可換候補の $h=1$ の方程式は
$G_{\alpha^{(1)}}+100G_{\alpha^{(2)}}$ であり、
定理\ref{cmo:thm:commutative-discriminant-candidate}の $c=0$ の滑らかな五様体を定める。
射影族の滑らかさの開性により、$R_1$ 上の候補は smooth proper である。
非可換側は定理\ref{mix:thm:h-semilocal-smooth}による。
これらの性質は $\widetilde R$ および $R'$ への底変更後も成立する。

まず、一般点の核を残余圏の inclusion と projection に合成し、
全 coherent 圏の間の核とする。
体 $K'$ 上での非可換射影積の核の定理
\cite[Theorem~4.15]{CMOBallardFarman}により、
これは混合二重次数射影積の perfect 対象である。
smooth proper 性により有界 coherent 複体で表示できる。
さらに、有限個の射と $d^2=0$ の関係を十分高い二重次数で持ち上げれば、
有限生成二重次数加群の有界複体で表示できる。
ここでは商圏の射が高次数の切断上の射として表示されることと、
有限生成 torsion 加群の十分高い二重次数が零であることを使う。

この有界複体の各項の有限個の斉次生成元を選ぶ。
最小の複体次数から順に、それらと前項の微分の像が生成する
$\widetilde R$ 上の graded lattice を取る。
各 lattice は一般点の項の部分加群なので $\widetilde R$-torsionfree であり、
各内部次数は有限自由である。
微分の像を加えたので lattice は有界複体をなし、$d^2=0$ はそのまま保たれる。
各項の一般ファイバーは元の項全体である。
これにより、一般点核を延長する通常の coherent 複体
$\widetilde P_{\widetilde R}$ を得る。

相対混合積では、各 coherent lattice の特殊ファイバーは compact である。
補題\ref{mix:lem:coherent-compact-lift}により各項が compact となるから、
有界複体 $\widetilde P_{\widetilde R}$ も perfect である。
両側の四つの例外線束による残余射影を合成する。
それらの射影核は有限個の evaluation cone からなる perfect 核である。
得られる残余双加群を $P_{\widetilde R}$ とすれば、
一般点では既に残余圏に属する核を射影しているので、
主張した一般点との同定を満たす。
各閉点で局所化すれば式\eqref{mix:eq:h-integral-kernel}を得る。
非 henselian な $R_1$ 上では $R'$ 自体が有限 $R_1$-加群とは限らないため、
有限性は局所化する前の $\widetilde R/R_1$ に対して主張している。
相対 smooth proper 性からこの核の左右随伴も perfect 核で表される。
\end{proof}

\begin{remark}[選択済み一般点核との相対整合性]
\label{mix:rem:h-integral-defect}
定理\ref{mix:thm:h-integral-kernel}は整核の存在を与えるが、その選択は一意ではない。
核から構成した随伴の unit と counit の cone は、
$R'$-perfect で、一般ファイバーが零なので、有限長の cohomology を持つ。
命題\paperref{main}{grl:prop:trait-kernel-test}の非負整数
$d_{R'}(P_{R'})$ をこの二つの cone の cohomology の長さの総和とすると、
\[
 \Phi_{P_{R'}\otimes_{R'}\C}\text{ が同値}
 \quad\Longleftrightarrow\quad d_{R'}(P_{R'})=0
\]
である。
有限 DVR 拡大ではこの defect は分岐指数倍となるので、
既に正の defect を持つ延長を底変更するだけでは消えない。
$h=1$ で何らかの比較核が存在することは、
定理\ref{cmo:thm:fixed-zero-commutative-kernel}により別に証明されている。
その直接構成の核が、ここで固定した $P_{K'}$ の特殊化として得られることは
まだ主張していない。
残る相対問題は、指定した一般点核を延長する整核の中から
$d_{R'}(P_{R'})=0$ となるものを構成することである。
固定ファイバーでの同値の存在、相対 smooth proper 模型の存在、
選択済み一般点核の可逆な整延長は、以上のように異なる主張である。
\end{remark}
